
This chapter is the principal new contribution of the major evaluation. It addresses the question raised at the minor evaluation: \emph{``If you cache state, how do you prevent two users from booking the same single piece?''}

\section{Problem Statement}
Let $r$ be a single, indivisible resource: the last unit of an inventory item, a unique reservation slot, or a per-second gateway token. Let $R_1, R_2, \ldots, R_k$ be concurrent client requests, each of which would consume $r$. Exactly \emph{one} request must succeed and consume $r$; all others must observe a definitive ``unavailable'' response. The system must satisfy:
\begin{itemize}[leftmargin=*]
    \item \textbf{Safety:} $r$ is never allocated more than once.
    \item \textbf{Liveness:} If $r$ is available and at least one request is well-formed, some request succeeds in bounded time.
    \item \textbf{Latency budget:} The contention resolution adds at most 5--10\,ms to the p95 of the protected operation.
    \item \textbf{Observability:} Every contention outcome is auditable.
\end{itemize}

The challenge is that the cached path designed in Chapter~4 is inherently \emph{stale-tolerant} --- which is precisely what makes it unsafe for single-resource state. We therefore route single-resource operations through a \emph{contention layer} that bypasses or extends the cache as required.

\section{Threat Model: How Caches Cause Double-Allocation}
Consider the naive implementation: a cached value \texttt{stock = 1} is read by two concurrent requests $R_1$ and $R_2$. Both decide they may proceed, both issue a \texttt{decrement} on the backend, and both succeed --- the backend now reports \texttt{stock = -1}, and one customer will receive a unit that does not exist. This is the canonical \emph{lost-update / overselling} race~\cite{oneuptime2026inventory, amitavroy2026hotel}. Figure~\ref{fig:race} sketches the timeline. The contention layer must prevent every variant of this race.

\begin{figure}[h]
\centering
\resizebox{0.9\textwidth}{!}{%
\begin{tikzpicture}[
    >=Stealth, thick,
    every node/.style={font=\footnotesize}
]
    \draw[->, gray] (0,2) -- (12,2) node[right]{time};
    \draw[->, gray] (0,0) -- (12,0) node[right]{time};
    \node[left] at (0,2) {$R_1$};
    \node[left] at (0,0) {$R_2$};

    \node[draw, fill=blue!10, rounded corners, minimum width=1.6cm] (r1read) at (1.6,2) {GET stock};
    \node[draw, fill=blue!10, rounded corners, minimum width=1.6cm] (r2read) at (3.0,0) {GET stock};
    \node[font=\scriptsize, above=0.05cm of r1read] {returns 1};
    \node[font=\scriptsize, below=0.05cm of r2read] {returns 1};

    \node[draw, fill=yellow!20, rounded corners, minimum width=1.6cm] (r1chk) at (4.6,2) {check $\geq$ 1};
    \node[draw, fill=yellow!20, rounded corners, minimum width=1.6cm] (r2chk) at (5.4,0) {check $\geq$ 1};

    \node[draw, fill=red!15, rounded corners, minimum width=1.8cm] (r1dec) at (7.6,2) {DECR stock};
    \node[draw, fill=red!15, rounded corners, minimum width=1.8cm] (r2dec) at (8.6,0) {DECR stock};
    \node[font=\scriptsize, above=0.05cm of r1dec] {ok, $=0$};
    \node[font=\scriptsize, below=0.05cm of r2dec] {ok, $=-1$ \textbf{!!}};

    \node[font=\scriptsize, draw, dashed, fill=red!5, rounded corners, align=center] at (10.6,1) {Both succeed.\\Stock oversold.};
\end{tikzpicture}%
}
\caption{Lost-update race on a single-resource decrement. $R_1$ and $R_2$ each read \texttt{stock=1} from the cache (or from Redis without an atomic guard), each independently passes the availability check, and each issues a successful decrement --- leaving the backend in an inconsistent \texttt{stock=-1} state.}
\label{fig:race}
\end{figure}

\section{Technique 1: Distributed Locking with Fencing Tokens}
\label{sec:lock}

\subsection{Mechanism}
A distributed lock --- typically Redlock~\cite{kleppmann2016locking, anonymous2025distlock} on Redis or a lease in etcd --- is acquired on a per-resource key (e.g., \texttt{lock:inventory:sku-42}) before the read-modify-write sequence. Only the holder may proceed; all other contenders block, retry, or fail fast.

\subsection{Fencing Tokens}
Plain locks are vulnerable to GC pauses, network partitions, and clock drift: a lock holder may believe it still holds the lock when another client has already taken over. The standard mitigation is a \emph{monotonically increasing fencing token} returned with the lock; the storage backend must reject writes carrying a token lower than the highest seen.

\begin{lstlisting}[language=Haskell, caption=Distributed lock acquisition with fencing token (Haskell)]
acquireLock :: Connection -> ByteString -> NominalDiffTime -> IO (Maybe FenceToken)
acquireLock conn key ttl = do
  token <- nextFenceToken conn key
  result <- runRedis conn $
    set key (encode token) [Ex (round ttl), Nx]
  case result of
    Right Ok -> pure (Just token)
    _        -> pure Nothing

withLock :: Connection -> ByteString -> NominalDiffTime
         -> (FenceToken -> IO a) -> IO (Maybe a)
withLock conn key ttl action = do
  mtok <- acquireLock conn key ttl
  case mtok of
    Nothing  -> pure Nothing
    Just tok -> Just <$> action tok `finally` releaseLock conn key tok
\end{lstlisting}

\subsection{Trade-offs}
Locks add 1--3\,ms of acquisition latency and serialise contended operations, capping per-resource throughput at $1 / (\text{critical-section duration})$. They are appropriate for \emph{coarse-grained, low-contention} resources (e.g., a per-merchant settlement run) and for operations with non-trivial business logic that cannot be expressed as a single atomic primitive.

\section{Technique 2: Atomic Decrement via Redis Lua Scripts}
\label{sec:lua}

\subsection{Mechanism}
For inventory-style operations, the entire ``check stock $\geq$ 1; decrement; return'' sequence is packaged in a Lua script. Redis executes Lua scripts atomically (single-threaded server model), eliminating the read-modify-write race without an explicit lock~\cite{oneuptime2026lua, dev2025redislua, oneuptime2026inventory}.

\begin{lstlisting}[language={[5.2]Lua}, caption=Atomic conditional decrement (Lua / Redis)]
-- KEYS[1] = inventory key, ARGV[1] = quantity requested
local available = tonumber(redis.call('GET', KEYS[1]) or '0')
local want      = tonumber(ARGV[1])
if available >= want then
    redis.call('DECRBY', KEYS[1], want)
    return {1, available - want}        -- {ok, remaining}
else
    return {0, available}               -- {failed, current}
end
\end{lstlisting}

\subsection{Integration with the Cache}
The cached path may still serve \emph{display reads} of the stock count (with explicit ``approximate'' semantics for the UI), but every \emph{authoritative} decrement bypasses the cache and goes through the Lua script. After a successful decrement, the script's return value is published on a Redis Pub/Sub channel; the server-local cache subscribes and invalidates its entry, bringing the cached display back in line within tens of milliseconds.

\subsection{Trade-offs}
Atomic Lua delivers the highest throughput per resource (Redis can execute well over $10^5$ short scripts per second per shard) and the lowest latency overhead ($<$1\,ms). Its downside is expressiveness: the entire critical section must fit in a Lua script that does not call out to other services. For business logic that requires external calls, fall back to Technique~1.

\section{Technique 3: Optimistic Concurrency Control with Versioned Cache Entries}

\subsection{Mechanism}
Each cache entry carries a monotonically increasing \emph{version number}. A reader records the version it observed; before committing a decrement, it issues a \texttt{CAS}-style operation (\texttt{WATCH}/\texttt{MULTI}/\texttt{EXEC} in Redis~\cite{redis2024inventory} or a conditional update in PostgreSQL using \texttt{WHERE version = ?}) that succeeds only if the version is unchanged. On conflict, the operation retries.

\begin{lstlisting}[language=SQL, caption=Versioned conditional update (PostgreSQL)]
-- Read
SELECT stock, version FROM inventory WHERE sku = $1;

-- Write (must succeed only if version unchanged)
UPDATE inventory
   SET stock   = stock - 1,
       version = version + 1
 WHERE sku     = $1
   AND version = $2
   AND stock   >= 1
RETURNING stock, version;
\end{lstlisting}

\subsection{Trade-offs}
OCC is ideal for \emph{low-contention} workloads where conflicts are rare; under high contention the retry rate climbs and effective throughput collapses. It avoids the head-of-line blocking of locks but pays a cost in wasted work~\cite{wang2024optiql, zhou2025ccaas}. We use OCC for merchant-configuration writes from the control plane (intrinsically rare events) and \emph{not} for inventory.

\section{Technique 4: PN-Counter CRDTs for Eventually-Consistent Reservations}

\subsection{Mechanism}
A PN-Counter~\cite{shapiro2011crdt, duncan2025crdt} is a CRDT consisting of two grow-only counters per replica: $P_i$ (positive) and $N_i$ (negative). The value is $\sum_i P_i - \sum_i N_i$, and replicas merge by element-wise maximum on each component. Decrements are recorded locally and propagated asynchronously; replicas converge under any communication pattern.

\subsection{Application: Soft Reservations}
PN-Counters are appropriate when an \emph{over-allocation tolerance} exists at the business level --- e.g., a flash sale advertising ``$\sim$1{,}000 units'' where a small overshoot is acceptable, or a soft-hold for cart-abandonment analytics. They are \emph{not} appropriate for hard inventory limits where every unit is precious.

\subsection{Trade-offs}
CRDTs eliminate coordination latency entirely (writes are local, never blocked) and survive partitions trivially. The cost is the loss of a hard upper bound: under sufficiently bursty contention they can over-decrement~\cite{ably2024crdt, ijsoc2025inventory}. We classify CRDTs as a \emph{relaxed} contention primitive and gate their use behind explicit business-class flags.

\section{Technique 5: Single-Flight Request Coalescing}
\label{sec:singleflight}

\subsection{Mechanism}
Single-flight (also called request coalescing) ensures that concurrent requests for the same key share a \emph{single} backend call. The first request acquires an in-flight slot and issues the call; all later requests for the same key during the window block on the slot's promise and receive the same response~\cite{1xapi2026singleflight, oneuptime2026stampede, oneuptime2026coalescing}.

\begin{lstlisting}[language=Haskell, caption={Single-flight wrapper around a backend fetch (Haskell, STM)}]
data SingleFlight k v = SingleFlight (TVar (HashMap k (TMVar (Either SomeException v))))

singleFlight :: (Eq k, Hashable k)
             => SingleFlight k v -> k -> IO v -> IO v
singleFlight (SingleFlight ref) k action = do
  (slot, owner) <- atomically $ do
    m <- readTVar ref
    case HM.lookup k m of
      Just s  -> pure (s, False)        -- follower
      Nothing -> do
        s <- newEmptyTMVar
        writeTVar ref (HM.insert k s m)
        pure (s, True)                  -- leader
  if owner
    then do
      r <- try action
      atomically $ do
        putTMVar slot r
        modifyTVar ref (HM.delete k)
      either throwIO pure r
    else atomically (readTMVar slot) >>= either throwIO pure
\end{lstlisting}

\subsection{Why Single-Flight Helps the ``Single Piece'' Question}
Even if the underlying backend uses any of Techniques 1--4, single-flight collapses the \emph{front of the contention curve}: instead of $k$ identical decrement attempts arriving at the backend simultaneously, exactly one arrives, and the other $k-1$ followers receive the leader's result. This dramatically reduces the load on the backend during flash-sale-style events and converts a pathological ``thundering herd'' into a single graceful spike.

It is important to note, however, that single-flight is \emph{not} a replacement for atomicity --- it merely deduplicates simultaneous attempts on the same node. Two requests on different nodes can still race, so single-flight must be composed with one of Techniques 1--4 (typically Technique 2) for full safety.

\section{Hybrid Architecture}

Figure~\ref{fig:contention} illustrates how a request flows through the integrated two-layer system.

\begin{figure}[h]
\centering
\resizebox{\textwidth}{!}{%
\begin{tikzpicture}[
    node distance=0.7cm and 0.9cm,
    box/.style={rectangle, draw, rounded corners, minimum height=0.85cm, minimum width=2.4cm, align=center, font=\small, fill=blue!5},
    decision/.style={diamond, draw, aspect=2, inner sep=1pt, align=center, font=\scriptsize, fill=yellow!15},
    primitive/.style={rectangle, draw, minimum height=0.85cm, minimum width=2.3cm, align=center, font=\scriptsize, fill=green!10},
    arrow/.style={-Stealth, thick}
]
    \node[box] (req) {Incoming\\Request};
    \node[box, right=of req] (t1) {Tier-1\\Thread Cache};
    \node[decision, right=of t1] (op) {Read or\\write?};

    \node[box, above right=0.3cm and 0.8cm of op] (sf) {Single-Flight\\Coalescer};
    \node[box, right=of sf] (t2) {Tier-2\\Server Cache};
    \node[box, right=of t2] (redis) {Redis /\\Backend};

    \node[primitive, below=1.2cm of sf] (lua) {Lua atomic\\decrement};
    \node[primitive, right=of lua] (lock) {Lock + fence\\token};
    \node[primitive, right=of lock] (occ) {OCC\\(versioned)};

    \draw[arrow] (req) -- (t1);
    \draw[arrow] (t1) -- (op);
    \draw[arrow] (op.north) |- node[pos=0.25, left, font=\scriptsize]{read} (sf.west);
    \draw[arrow] (op.south) |- node[pos=0.25, left, font=\scriptsize]{write} (lua.west);
    \draw[arrow] (sf) -- (t2);
    \draw[arrow] (t2) -- (redis);
    \draw[arrow] (lua.east) -- (lock.west);
    \draw[arrow] (lock.east) -- (occ.west);
    \draw[arrow] (occ.north) -- (redis.south);
    \draw[arrow, dashed] (redis.north) to[bend left=30] node[above, font=\scriptsize]{invalidate (Pub/Sub)} (t2.north);
\end{tikzpicture}%
}
\caption{Integrated two-layer architecture. Reads traverse Tier-1, the single-flight coalescer, and Tier-2, reaching Redis only on a Tier-2 miss. Authoritative writes bypass the cache and select a contention primitive (Lua, lock+fence, or OCC). Successful writes publish invalidation events that propagate back to Tier-2.}
\label{fig:contention}
\end{figure}

\subsection{Decision Matrix}
Table~\ref{tab:matrix} maps each data class to its cache tier and contention primitive.

\begin{table}[h]
\centering
\footnotesize
\setlength{\tabcolsep}{4pt}
\renewcommand{\arraystretch}{1.15}
\begin{tabular}{p{3.4cm}p{1.7cm}p{2.5cm}p{2.0cm}p{2.6cm}}
\toprule
\textbf{Data class} & \textbf{Cache?} & \textbf{Primary} & \textbf{Coalescing} & \textbf{Notes} \\
\midrule
Merchant config (read)   & Yes (Tier-2) & ---              & Single-flight & Cache-eligible \\
Routing table (read)     & Yes (Tier-2) & ---              & Single-flight & Cache-eligible \\
Feature flag (read)      & Yes (Tier-2) & ---              & Single-flight & Cache-eligible \\
Inventory (read display) & Yes (Tier-2) & ---              & Single-flight & Marked approximate \\
Inventory (decrement)    & No           & Lua atomic       & Single-flight & Hard limit \\
Account balance debit    & No           & Distributed lock & Single-flight & Coarse, fence token \\
Reservation slot         & No           & Lua atomic       & Single-flight & Per-slot key \\
Soft hold (cart)         & Tier-2       & PN-Counter       & ---           & Over-allocation OK \\
Merchant config (write)  & Invalidate   & OCC versioned    & ---           & Low contention \\
\bottomrule
\end{tabular}
\caption{Decision matrix mapping data classes to cache tier and primary contention primitive.}
\label{tab:matrix}
\end{table}

\noindent The hybrid architecture composes the primitives as follows:
\begin{enumerate}[leftmargin=*]
    \item \textbf{Reads} go through the two-tier cache, fronted by single-flight at both tiers.
    \item \textbf{Authoritative writes on single resources} bypass the cache and route to the appropriate primary primitive (Lua, lock, OCC, or PN-Counter) according to Table~\ref{tab:matrix}.
    \item \textbf{Cache invalidation} on successful writes is published via Redis Pub/Sub; subscribers invalidate their server-local entries within tens of milliseconds.
    \item \textbf{Fence tokens} are stamped on every authoritative write that uses Technique~1, and the storage backend rejects writes with stale tokens.
\end{enumerate}

\section{Comparative Analysis}

\begin{table}[h]
\centering
\footnotesize
\setlength{\tabcolsep}{4pt}
\renewcommand{\arraystretch}{1.2}
\resizebox{\textwidth}{!}{%
\begin{tabular}{p{3.4cm}p{2.0cm}p{2.4cm}p{2.6cm}p{1.8cm}p{2.2cm}}
\toprule
\textbf{Property} & \textbf{Lock} & \textbf{Lua} & \textbf{OCC} & \textbf{CRDT} & \textbf{Single-flight} \\
\midrule
Hard upper bound          & Yes              & Yes                  & Yes                    & No        & N/A \\
Per-key throughput        & Low              & Very high            & Medium--high           & Very high & N/A \\
Latency overhead          & 1--3\,ms         & $<$1\,ms             & $<$1\,ms (no conflict) & 0\,ms     & 0--1\,ms \\
Behaviour under partition & Blocks / fails   & Fails fast           & Fails fast             & Continues & Continues \\
Expressiveness            & High (any logic) & Low (Lua only)       & Medium (SQL / CAS)     & Counter only & N/A \\
Composes with cache       & Bypass           & Bypass + invalidate  & Versioned cache        & Cached    & Front of cache \\
\bottomrule
\end{tabular}%
}
\caption{Comparative properties of the five contention primitives.}
\label{tab:compare}
\end{table}

% =============================================================================
